\IfFileExists{stacks-project.cls}{%
\documentclass{stacks-project}
}{%
\documentclass{amsart}
}

% For dealing with references we use the comment environment
\usepackage{verbatim}
\newenvironment{reference}{\comment}{\endcomment}
%\newenvironment{reference}{}{}
\newenvironment{slogan}{\comment}{\endcomment}
\newenvironment{history}{\comment}{\endcomment}

% For commutative diagrams we use Xy-pic
\usepackage[all]{xy}

% We use 2cell for 2-commutative diagrams.
\xyoption{2cell}
\UseAllTwocells

% We use multicol for the list of chapters between chapters
\usepackage{multicol}

% This is generally recommended for better output
\usepackage{lmodern}
\usepackage[T1]{fontenc}

% For cross-file-references
\usepackage{xr-hyper}

% Package for hypertext links:
\usepackage{hyperref}

% For any local file, say "hello.tex" you want to link to please
% use \externaldocument[hello-]{hello}
\externaldocument[introduction-]{introduction}
\externaldocument[conventions-]{conventions}
\externaldocument[sets-]{sets}
\externaldocument[categories-]{categories}
\externaldocument[topology-]{topology}
\externaldocument[sheaves-]{sheaves}
\externaldocument[sites-]{sites}
\externaldocument[stacks-]{stacks}
\externaldocument[fields-]{fields}
\externaldocument[algebra-]{algebra}
\externaldocument[brauer-]{brauer}
\externaldocument[homology-]{homology}
\externaldocument[derived-]{derived}
\externaldocument[simplicial-]{simplicial}
\externaldocument[more-algebra-]{more-algebra}
\externaldocument[smoothing-]{smoothing}
\externaldocument[modules-]{modules}
\externaldocument[sites-modules-]{sites-modules}
\externaldocument[injectives-]{injectives}
\externaldocument[cohomology-]{cohomology}
\externaldocument[sites-cohomology-]{sites-cohomology}
\externaldocument[dga-]{dga}
\externaldocument[dpa-]{dpa}
\externaldocument[sdga-]{sdga}
\externaldocument[hypercovering-]{hypercovering}
\externaldocument[schemes-]{schemes}
\externaldocument[constructions-]{constructions}
\externaldocument[properties-]{properties}
\externaldocument[morphisms-]{morphisms}
\externaldocument[coherent-]{coherent}
\externaldocument[divisors-]{divisors}
\externaldocument[limits-]{limits}
\externaldocument[varieties-]{varieties}
\externaldocument[topologies-]{topologies}
\externaldocument[descent-]{descent}
\externaldocument[perfect-]{perfect}
\externaldocument[more-morphisms-]{more-morphisms}
\externaldocument[flat-]{flat}
\externaldocument[groupoids-]{groupoids}
\externaldocument[more-groupoids-]{more-groupoids}
\externaldocument[etale-]{etale}
\externaldocument[chow-]{chow}
\externaldocument[intersection-]{intersection}
\externaldocument[pic-]{pic}
\externaldocument[weil-]{weil}
\externaldocument[adequate-]{adequate}
\externaldocument[dualizing-]{dualizing}
\externaldocument[duality-]{duality}
\externaldocument[discriminant-]{discriminant}
\externaldocument[derham-]{derham}
\externaldocument[local-cohomology-]{local-cohomology}
\externaldocument[algebraization-]{algebraization}
\externaldocument[curves-]{curves}
\externaldocument[resolve-]{resolve}
\externaldocument[models-]{models}
\externaldocument[functors-]{functors}
\externaldocument[equiv-]{equiv}
\externaldocument[pione-]{pione}
\externaldocument[etale-cohomology-]{etale-cohomology}
\externaldocument[proetale-]{proetale}
\externaldocument[relative-cycles-]{relative-cycles}
\externaldocument[more-etale-]{more-etale}
\externaldocument[trace-]{trace}
\externaldocument[crystalline-]{crystalline}
\externaldocument[spaces-]{spaces}
\externaldocument[spaces-properties-]{spaces-properties}
\externaldocument[spaces-morphisms-]{spaces-morphisms}
\externaldocument[decent-spaces-]{decent-spaces}
\externaldocument[spaces-cohomology-]{spaces-cohomology}
\externaldocument[spaces-limits-]{spaces-limits}
\externaldocument[spaces-divisors-]{spaces-divisors}
\externaldocument[spaces-over-fields-]{spaces-over-fields}
\externaldocument[spaces-topologies-]{spaces-topologies}
\externaldocument[spaces-descent-]{spaces-descent}
\externaldocument[spaces-perfect-]{spaces-perfect}
\externaldocument[spaces-more-morphisms-]{spaces-more-morphisms}
\externaldocument[spaces-flat-]{spaces-flat}
\externaldocument[spaces-groupoids-]{spaces-groupoids}
\externaldocument[spaces-more-groupoids-]{spaces-more-groupoids}
\externaldocument[bootstrap-]{bootstrap}
\externaldocument[spaces-pushouts-]{spaces-pushouts}
\externaldocument[spaces-chow-]{spaces-chow}
\externaldocument[groupoids-quotients-]{groupoids-quotients}
\externaldocument[spaces-more-cohomology-]{spaces-more-cohomology}
\externaldocument[spaces-simplicial-]{spaces-simplicial}
\externaldocument[spaces-duality-]{spaces-duality}
\externaldocument[formal-spaces-]{formal-spaces}
\externaldocument[restricted-]{restricted}
\externaldocument[spaces-resolve-]{spaces-resolve}
\externaldocument[formal-defos-]{formal-defos}
\externaldocument[defos-]{defos}
\externaldocument[cotangent-]{cotangent}
\externaldocument[examples-defos-]{examples-defos}
\externaldocument[algebraic-]{algebraic}
\externaldocument[examples-stacks-]{examples-stacks}
\externaldocument[stacks-sheaves-]{stacks-sheaves}
\externaldocument[criteria-]{criteria}
\externaldocument[artin-]{artin}
\externaldocument[quot-]{quot}
\externaldocument[stacks-properties-]{stacks-properties}
\externaldocument[stacks-morphisms-]{stacks-morphisms}
\externaldocument[stacks-limits-]{stacks-limits}
\externaldocument[stacks-cohomology-]{stacks-cohomology}
\externaldocument[stacks-perfect-]{stacks-perfect}
\externaldocument[stacks-introduction-]{stacks-introduction}
\externaldocument[stacks-more-morphisms-]{stacks-more-morphisms}
\externaldocument[stacks-geometry-]{stacks-geometry}
\externaldocument[moduli-]{moduli}
\externaldocument[moduli-curves-]{moduli-curves}
\externaldocument[examples-]{examples}
\externaldocument[exercises-]{exercises}
\externaldocument[guide-]{guide}
\externaldocument[desirables-]{desirables}
\externaldocument[coding-]{coding}
\externaldocument[obsolete-]{obsolete}
\externaldocument[fdl-]{fdl}
\externaldocument[index-]{index}

% Theorem environments.
%
\theoremstyle{plain}
\newtheorem{theorem}[subsection]{Theorem}
\newtheorem{proposition}[subsection]{Proposition}
\newtheorem{lemma}[subsection]{Lemma}

\theoremstyle{definition}
\newtheorem{definition}[subsection]{Definition}
\newtheorem{example}[subsection]{Example}
\newtheorem{exercise}[subsection]{Exercise}
\newtheorem{situation}[subsection]{Situation}

\theoremstyle{remark}
\newtheorem{remark}[subsection]{Remark}
\newtheorem{remarks}[subsection]{Remarks}

\numberwithin{equation}{subsection}

% Macros
%
\def\lim{\mathop{\mathrm{lim}}\nolimits}
\def\colim{\mathop{\mathrm{colim}}\nolimits}
\def\Spec{\mathop{\mathrm{Spec}}}
\def\Hom{\mathop{\mathrm{Hom}}\nolimits}
\def\Ext{\mathop{\mathrm{Ext}}\nolimits}
\def\SheafHom{\mathop{\mathcal{H}\!\mathit{om}}\nolimits}
\def\SheafExt{\mathop{\mathcal{E}\!\mathit{xt}}\nolimits}
\def\Sch{\mathit{Sch}}
\def\Mor{\mathop{\mathrm{Mor}}\nolimits}
\def\Ob{\mathop{\mathrm{Ob}}\nolimits}
\def\Sh{\mathop{\mathit{Sh}}\nolimits}
\def\NL{\mathop{N\!L}\nolimits}
\def\CH{\mathop{\mathrm{CH}}\nolimits}
\def\proetale{{pro\text{-}\acute{e}tale}}
\def\etale{{\acute{e}tale}}
\def\QCoh{\mathit{QCoh}}
\def\Ker{\mathop{\mathrm{Ker}}}
\def\Im{\mathop{\mathrm{Im}}}
\def\Coker{\mathop{\mathrm{Coker}}}
\def\Coim{\mathop{\mathrm{Coim}}}

% Boxtimes
%
\DeclareMathSymbol{\boxtimes}{\mathbin}{AMSa}{"02}

%
% Macros for moduli stacks/spaces
%
\def\QCohstack{\mathcal{QC}\!\mathit{oh}}
\def\Cohstack{\mathcal{C}\!\mathit{oh}}
\def\Spacesstack{\mathcal{S}\!\mathit{paces}}
\def\Quotfunctor{\mathrm{Quot}}
\def\Hilbfunctor{\mathrm{Hilb}}
\def\Curvesstack{\mathcal{C}\!\mathit{urves}}
\def\Polarizedstack{\mathcal{P}\!\mathit{olarized}}
\def\Complexesstack{\mathcal{C}\!\mathit{omplexes}}
% \Pic is the operator that assigns to X its picard group, usage \Pic(X)
% \Picardstack_{X/B} denotes the Picard stack of X over B
% \Picardfunctor_{X/B} denotes the Picard functor of X over B
\def\Pic{\mathop{\mathrm{Pic}}\nolimits}
\def\Picardstack{\mathcal{P}\!\mathit{ic}}
\def\Picardfunctor{\mathrm{Pic}}
\def\Deformationcategory{\mathcal{D}\!\mathit{ef}}
\title{A map from a compact object into a nicely filtered dg module factors through a finitely filtered submodule}
\author{The Stacks Project Authors}
\date{stacks-project:dga.tex:local-snapshot-20261012; GNU FDL 1.2 or later}
\begin{document}
\maketitle
\noindent\textit{Editorial note.} The counted statement is the author's own lemma environment labelled lemma-factor-through-nicer (source lines 6433-6451); the proof environment begins two lines below it and is its own proof, closing at source line 6495.\par\medskip
\noindent\textit{Editorial note.} No \textbackslash{}ref or \textbackslash{}cite occurs anywhere inside the retained span, so the package declares neither a source pointer nor a bibliography entry; the compactness and filtration notions are part of the author's surrounding development in the same chapter and are named rather than cited by label.\par\medskip
\noindent\textit{Editorial note.} The proof is an induction whose hypothesis is stated inside the proof itself; the three listed clauses of that hypothesis are the author's own wording and are reproduced verbatim.\par\medskip
\noindent\textit{Editorial note.} The retained preamble is the chapter's own minimal preamble (sources/stacks-dga/preamble.tex, 217 lines), a verbatim excerpt of the native Xy-pic setup, the native theorem declarations with the mathematical macros, and the project's own no-op reference environments; no macro was invented or redefined.\par\medskip
\noindent\textit{Editorial scope.} One original result of the Stacks Project's Differential Graded Algebra chapter, counted once: a map from a compact object into a nicely filtered dg module factors through a submodule with finite direct-sum quotients.\par\medskip
\noindent\textit{Allowed standard prerequisites (editorial).}\begin{itemize}
\item Compact object of a triangulated category. An object E is compact when the functor it represents commutes with arbitrary direct sums, so that any morphism from E into a direct sum factors through a finite subsum.. It is invoked at source lines 6458-6462 to enlarge the finite index set so that E -> P factors through the constructed submodule.
\item Finite filtration of a dg module by shifted copies of A. The module P carries a finite increasing filtration by dg submodules whose successive quotients are direct sums of A shifted by integers.. The induction runs down the filtration levels and the quotient description is what is preserved at every step (source lines 6436-6440, 6465-6470).
\item Homogeneous generators and the differential on a quotient piece. Each graded piece is free on homogeneous generators of known degrees, and the differential of a generator decomposes into a combination of the other generators plus a term in the next lower filtration level.. The choice of the x\_\{i,j\} and the displayed formula for d(x\_\{i,j\}) drive the construction (source lines 6456-6465).
\end{itemize}\medskip
\begin{lemma}
\label{lemma-factor-through-nicer}
Let $(A, \text{d})$ be a differential graded algebra. Let $E$ be a compact
object of $D(A, \text{d})$. Let $P$ be a differential graded $A$-module
which has a finite filtration
$$
0 = F_{-1}P \subset F_0P \subset F_1P \subset \ldots \subset F_nP = P
$$
by differential graded submodules such that
$$
F_{i + 1}P/F_iP \cong \bigoplus\nolimits_{j \in J_i} A[k_{i, j}]
$$
as differential graded $A$-modules for some sets $J_i$ and integers $k_{i, j}$.
Let $E \to P$ be a morphism of $D(A, \text{d})$.
Then there exists a differential graded submodule $P' \subset P$ such that
$F_{i + 1}P \cap P'/(F_iP \cap P')$ is equal to
$\bigoplus_{j \in J'_i} A[k_{i, j}]$ for some finite subsets
$J'_i \subset J_i$ and such that $E \to P$ factors through $P'$.
\end{lemma}

\begin{proof}
We will prove by induction on $-1 \leq m \leq n$ that there exists
a differential graded submodule $P' \subset P$  such that
\begin{enumerate}
\item $F_mP \subset P'$,
\item for $i \geq m$ the quotient $F_{i + 1}P \cap P'/(F_iP \cap P')$ is
isomorphic to $\bigoplus_{j \in J'_i} A[k_{i, j}]$ for some finite subsets
$J'_i \subset J_i$, and
\item $E \to P$ factors through $P'$.
\end{enumerate}
The base case is $m = n$ where we can take $P' = P$.

\medskip\noindent
Induction step. Assume $P'$ works for $m$.
For $i \geq m$ and $j \in J'_i$ let $x_{i, j} \in F_{i + 1}P \cap P'$
be a homogeneous element of degree $k_{i, j}$ whose image in
$F_{i + 1}P \cap P'/(F_iP \cap P')$ is the generator in
the summand corresponding to $j \in J_i$. The
$x_{i, j}$ generate $P'/F_mP$ as an $A$-module. Write
$$
\text{d}(x_{i, j}) = \sum x_{i', j'} a_{i, j}^{i', j'} + y_{i, j}
$$
with $y_{i, j} \in F_mP$ and $a_{i, j}^{i', j'} \in A$.
There exists a finite subset
$J'_{m - 1} \subset J_{m - 1}$ such that each $y_{i, j}$ maps to
an element of the submodule $\bigoplus_{j \in J'_{m - 1}} A[k_{m - 1, j}]$
of $F_mP/F_{m - 1}P$. Let $P'' \subset F_mP$ be the inverse
image of $\bigoplus_{j \in J'_{m - 1}} A[k_{m - 1, j}]$ under
the map $F_mP \to F_mP/F_{m - 1}P$. Then we see that the $A$-submodule
$$
P'' + \sum x_{i, j}A
$$
is a differential graded submodule of the type we are looking for. Moreover
$$
P'/(P'' + \sum x_{i, j}A) =
\bigoplus\nolimits_{j \in J_{m - 1} \setminus J'_{m - 1}} A[k_{m - 1, j}]
$$
Since $E$ is compact, the composition of the given map $E \to P'$
with the quotient map, factors through a finite direct subsum of
the module displayed above. Hence after enlarging $J'_{m - 1}$
we may assume $E \to P'$ factors through
$P'' + \sum x_{i, j}A$ as desired.
\end{proof}

\end{document}
